\section{指令周期、数据通路与流水线}

\term{冯·诺依曼式}{von Neumann}组织让指令和数据共享一个可寻址存储体系；
\term{Harvard 式}{Harvard}组织为指令和数据提供分离存储或接口。现代 CPU
常表现为混合结构：软件看到统一虚拟地址空间，内部却有分离的 L1 指令缓存和
数据缓存，更下层再汇合。

\begin{osgridlongtable}{L{0.19\linewidth}L{0.31\linewidth}L{0.40\linewidth}}
组织 & 优势 & 代价与现实变体 \\ \hline
\endfirsthead
组织 & 优势 & 代价与现实变体 \\ \hline
\endhead
统一存储 & 指令和数据可共享容量，程序易装载修改 & 取指与数据访问争用同一通路 \\ \hline
分离存储 & 可同时取指和访数据，权限/宽度可不同 & 两边容量固定，交换内容需要额外机制 \\ \hline
现代混合 & ISA 统一、近端缓存分离、下层汇合 & 必须维护指令/数据缓存与修改代码的一致性 \\ \hline
\end{osgridlongtable}

“程序也是 byte”表示指令可以存储和搬运，不表示 CPU 会把任意数据自动执行。
取指必须来自当前控制流选择的地址，并通过执行权限、模式和译码规则。

\topic{把 CPU 打开：数据通路与控制通路分别做什么}
\term{数据通路}{datapath}搬运并变换值，包括寄存器文件、ALU、多路选择器和
总线；\term{控制通路}{control path}根据指令、状态和条件产生选择、读写、
使能与下一状态。

\begin{center}
\begin{tikzpicture}[font=\footnotesize\sffamily,>=Latex]
  \node[os state,minimum width=23mm] (pc) at (0,4.6) {PC/RIP\\取指地址};
  \node[os hardware,minimum width=25mm] (memory) at (3.5,4.6)
    {指令存储接口\\返回机器码};
  \node[os state,minimum width=23mm] (ir) at (7.0,4.6)
    {指令寄存器 IR\\当前编码};
  \node[os block,minimum width=27mm] (decode) at (10.5,4.6)
    {译码/控制器\\产生控制信号};

  \node[os state,minimum width=28mm] (regs) at (0.5,1.5)
    {寄存器文件\\两个读口、一个写口};
  \node[os block,minimum width=28mm] (operands) at (4.0,1.5)
    {操作数选择\\MUX A / MUX B};
  \node[os block,minimum width=27mm] (alu) at (7.4,1.5)
    {ALU\\算术/地址/逻辑};
  \node[os state,minimum width=25mm] (result) at (10.5,1.5)
    {结果与标志};

  \draw[os bus] (pc) -- (memory);
  \draw[os bus] (memory) -- (ir);
  \draw[os bus] (ir) -- (decode);
  \draw[os bus] ([yshift=2mm]regs.east) -- ([yshift=2mm]operands.west);
  \draw[os bus] ([yshift=-2mm]regs.east) -- ([yshift=-2mm]operands.west);
  \draw[os bus] (operands) -- (alu);
  \draw[os bus] (alu) -- (result);

  \draw[os arrow,BookTeal] (decode.south) --
    node[right]{提交使能} (result.north);
  \draw[os arrow,BookTeal] ([xshift=-4mm]decode.south) -- ++(0,-0.55) -|
    node[pos=0.30,above]{选择与运算控制} (operands.north);

  \draw[os bus] (result.south) -- ++(0,-1.0) -|
    node[pos=0.45,below]{写回} (regs.south);
  \draw[os arrow] (result.east) -- ++(0.65,0) |- ++(0,4.2) -|
    node[pos=0.45,above]{顺序、分支或异常决定下一 RIP} (pc.north);
\end{tikzpicture}
\end{center}
\diagramnote{蓝色双线表示多 bit 数据，绿色细线表示控制选择。这是一张教学
数据通路；真实 x86 核心还包含乱序执行、预测、缓存等结构，几何位置也不同。}

\topic{寄存器文件为什么常有多个端口}
一条二操作数算术指令通常要同时读取两个源，再写一个目标。寄存器文件可提供
两个组合读端口和一个时钟写端口。地址译码选择寄存器，读多路选择器返回内容，
写使能只在提交时更新目标。

\begin{center}
\begin{tikzpicture}[font=\footnotesize\sffamily,>=Latex]
  \draw[rounded corners=3pt,draw=BookBlue,fill=BookCodeBg]
    (3,0) rectangle (8,4.8);
  \node[font=\bfseries] at (5.5,4.3) {寄存器文件};
  \foreach \y/\name in {3.4/R0,2.7/R1,2.0/R2,1.3/\(\cdots\),0.6/R15}
    \node[draw=BookRule,fill=white,minimum width=18mm,minimum height=5mm]
      at (5.5,\y) {\name};
  \node[os state] (ra) at (0,3.6) {read addr A};
  \node[os state] (rb) at (0,2.4) {read addr B};
  \node[os state] (wa) at (0,1.2) {write addr};
  \node[os state] (wd) at (5.5,-1.0) {write data + enable};
  \node[os state] (qa) at (11,3.3) {read data A};
  \node[os state] (qb) at (11,1.6) {read data B};
  \draw[os arrow] (ra) -- (3,3.6);
  \draw[os arrow] (rb) -- (3,2.4);
  \draw[os arrow] (wa) -- (3,1.2);
  \draw[os bus] (wd) -- (5.5,0);
  \draw[os bus] (8,3.3) -- (qa);
  \draw[os bus] (8,1.6) -- (qb);
\end{tikzpicture}
\end{center}

\topic{机器指令包含“做什么”和“对谁做”}
机器指令编码通常要表达部分或全部字段：

\begin{itemize}[leftmargin=2.2em]
  \item opcode：ADD、LOAD、JUMP 等操作；
  \item 源/目标寄存器编号；
  \item 立即数或位移；
  \item 寻址方式、操作数宽度和条件；
  \item 某些 ISA 中的权限或扩展前缀。
\end{itemize}

为建立直觉，定义一套\textbf{纯教学、不是 x86} 的 16-bit 定长指令。所有
指令都占 2 byte，最高 4 bit 是 opcode，其余字段随指令种类解释：

\begin{center}
\begin{tikzpicture}[font=\footnotesize\sffamily]
  \foreach \x/\w/\colorname/\title/\bits in {
    0/3.2/BookBlue/opcode/\(15..12\),
    3.2/3.2/BookTeal/字段 A/\(11..8\),
    6.4/6.4/BookAmber/字段 B/\(7..0\)} {
    \fill[\colorname!18] (\x,0) rectangle ++(\w,1.15);
    \draw[\colorname,thick] (\x,0) rectangle ++(\w,1.15);
    \node at (\x+\w/2,0.68) {\title};
    \node[below] at (\x+\w/2,0) {\bits};
  }
\end{tikzpicture}
\end{center}

\begin{osgridlongtable}{L{0.13\linewidth}L{0.19\linewidth}L{0.25\linewidth}L{0.33\linewidth}}
opcode & 汇编写法 & 字段解释 & 执行动作 \\ \hline
\endfirsthead
opcode & 汇编写法 & 字段解释 & 执行动作 \\ \hline
\endhead
\texttt{0x1} & \code{LOAD Rd, imm8} & A=Rd，B=8-bit 立即数 & \(R_d\leftarrow imm8\) \\ \hline
\texttt{0x2} & \code{ADD Rd, Rs} & A=Rd，B 高 4 bit=Rs & \(R_d\leftarrow R_d+R_s\) \\ \hline
\texttt{0x3} & \code{STORE [addr8], Rs} & A=Rs，B=byte 地址 & \(M[addr8]\leftarrow R_s\) \\ \hline
\texttt{0x4} & \code{JMP addr12} & 低 12 bit 合成目标 & \(PC\leftarrow addr12\) \\ \hline
\texttt{0x0} & \code{HALT} & 其余字段必须为 0 & 停在 Halt 状态 \\ \hline
\end{osgridlongtable}

\begin{workedexample}
\code{LOAD R0, 5} 的字段为 opcode=\texttt{1}、Rd=\texttt{0}、
imm8=\texttt{05}，16-bit 指令字是 \texttt{0x1005}。若 ROM 也采用
little-endian byte 排列，低地址先放 \texttt{05}，高地址再放
\texttt{10}。汇编写法、指令字和 ROM 中的 byte 顺序属于三个相邻层次。
\end{workedexample}

真实 x86 指令是可变长度编码，前缀、操作码、ModR/M、SIB、位移和立即数字段
会按指令变化。教学编码只用于把“哪几个 bit 控制哪条数据通路”看清楚。

\topic{让五条教学指令在 ROM 中跑一圈}
把程序从地址 \(\texttt{0x10}\) 开始排列：

\begin{center}
\begin{osgridtabular}{cL{0.29\linewidth}cL{0.29\linewidth}}
地址 & 汇编 & 指令字 & ROM 中两个 byte \\ \hline
\texttt{0x10} & \code{LOAD R0, 5} & \texttt{0x1005} & \texttt{05 10} \\ \hline
\texttt{0x12} & \code{LOAD R1, 3} & \texttt{0x1103} & \texttt{03 11} \\ \hline
\texttt{0x14} & \code{ADD R0, R1} & \texttt{0x2010} & \texttt{10 20} \\ \hline
\texttt{0x16} & \code{STORE [0x20], R0} & \texttt{0x3020} & \texttt{20 30} \\ \hline
\texttt{0x18} & \code{JMP 0x10} & \texttt{0x4010} & \texttt{10 40} \\ \hline
\end{osgridtabular}
\end{center}

复位先令 \(PC=\texttt{0x10}\)，R0、R1 和 RAM
\(\texttt{0x20}\) 均为 0。第一条 \code{LOAD} 可以拆成五拍：

\begin{osgridlongtable}{L{0.12\linewidth}L{0.20\linewidth}L{0.35\linewidth}L{0.23\linewidth}}
周期 & 控制状态 & 发生的动作 & 周期结束后的关键状态 \\ \hline
\endfirsthead
周期 & 控制状态 & 发生的动作 & 周期结束后的关键状态 \\ \hline
\endhead
1 & FetchAddr & PC 把 \texttt{0x10} 送到地址总线，发出读请求 & ROM 开始选择指令 byte \\ \hline
2 & FetchData & ROM 返回 \texttt{05 10}，IR 锁存 \texttt{0x1005} & \(PC=\texttt{0x12}\) \\ \hline
3 & Decode & 译码 opcode=\texttt{1}、Rd=R0、imm8=\texttt{5} & 控制器选择立即数输入 \\ \hline
4 & Execute & 数据通路形成候选值 5 & R0 仍为 0 \\ \hline
5 & Commit & R0 写使能在边沿生效 & \(R0=5\) \\ \hline
\end{osgridlongtable}

第二条 \code{LOAD} 以同样步骤得到 \(R1=3\)。第三条
\code{ADD R0, R1} 读取两个寄存器，由 ALU 算出 8，再把结果写回 R0。
第四条把 8 写到 RAM 地址 \(\texttt{0x20}\)。第五条不采用顺序
\(\texttt{0x1A}\)，而把 PC 改回 \(\texttt{0x10}\)。

\begin{osgridlongtable}{L{0.14\linewidth}L{0.25\linewidth}L{0.15\linewidth}L{0.15\linewidth}L{0.21\linewidth}}
指令完成后 & 刚完成的指令 & R0 & R1 & RAM[\texttt{0x20}] / PC \\ \hline
\endfirsthead
指令完成后 & 刚完成的指令 & R0 & R1 & RAM[\texttt{0x20}] / PC \\ \hline
\endhead
复位 & 尚未取指 & 0 & 0 & 0 / \texttt{0x10} \\ \hline
第 1 条 & \code{LOAD R0, 5} & 5 & 0 & 0 / \texttt{0x12} \\ \hline
第 2 条 & \code{LOAD R1, 3} & 5 & 3 & 0 / \texttt{0x14} \\ \hline
第 3 条 & \code{ADD R0, R1} & 8 & 3 & 0 / \texttt{0x16} \\ \hline
第 4 条 & \code{STORE [0x20], R0} & 8 & 3 & 8 / \texttt{0x18} \\ \hline
第 5 条 & \code{JMP 0x10} & 8 & 3 & 8 / \texttt{0x10} \\ \hline
\end{osgridlongtable}

\handsonexperiment
在纸上执行第二轮。注意第一条会把 R0 从 8 重新写成 5，所以每一轮写入
\(\texttt{0x20}\) 的仍是 8，而不是 16。只看最后一条 \code{JMP} 无法猜出
程序效果；必须按 PC、IR、寄存器和内存的状态变化逐条跟踪。

\topic{逐拍跟踪一条指令，才能知道每个部件何时有用}
设教学 CPU 每个状态占一个时钟周期：

\begin{center}
\begin{tikzpicture}[font=\footnotesize\sffamily,>=Latex]
  \node[os state] (fa) at (0,2.5) {FetchAddr};
  \node[os state] (fd) at (3.0,2.5) {FetchData};
  \node[os state] (decode) at (6.0,2.5) {Decode};
  \node[os state] (execute) at (9.0,2.5) {Execute};
  \node[os state,minimum width=27mm] (commit) at (12.0,2.5)
    {Commit\\完成后取下一条};
  \node[os state] (memory) at (9.0,0) {Memory};
  \node[os state,draw=BookRed,minimum width=27mm] (error) at (3.0,0)
    {Error/Halt\\保持错误现场};

  \draw[os arrow] (fa) -- (fd);
  \draw[os arrow] (fd) -- (decode);
  \draw[os arrow] (decode) -- (execute);
  \draw[os arrow] (execute) -- node[above]{ALU/JMP} (commit);
  \draw[os arrow] (execute) -- node[right]{LOAD/STORE} (memory);
  \draw[os arrow] (memory) -- node[above,sloped]{完成访存} (commit);
  \draw[os arrow,draw=BookRed] (fd) -- node[right]{取指错误} (error);
  \draw[os arrow,draw=BookRed] (memory) -- node[above]{数据错误} (error);
\end{tikzpicture}
\end{center}
\diagramnote{控制单元本身也是有限状态机。ALU/JMP 可从 Execute 直接进入
Commit，LOAD/STORE 先经过 Memory；取指或数据错误进入 Error/Halt。Commit
结束后回到 FetchAddr 的动作写在节点中，避免回边穿过错误路径。}

\begin{osgridlongtable}{L{0.13\linewidth}L{0.23\linewidth}L{0.31\linewidth}L{0.23\linewidth}}
拍 & 主状态 & 数据移动 & 控制动作 \\ \hline
\endfirsthead
拍 & 主状态 & 数据移动 & 控制动作 \\ \hline
\endhead
\(T_0\) & Fetch address & \(MAR\leftarrow PC\) & 发出 memory read \\ \hline
\(T_1\) & Fetch data & \(IR\leftarrow M[MAR]\) & 等待 ready，\(PC\leftarrow PC+2\) \\ \hline
\(T_2\) & Decode/read & \(A\leftarrow R[rd], B\leftarrow R[rs]\) & 译码 opcode、选择两个寄存器 \\ \hline
\(T_3\) & Execute & \(ALUOut\leftarrow A+B\) & 选择 ADD，形成标志候选 \\ \hline
\(T_4\) & Commit & \(R[rd]\leftarrow ALUOut\) & 写使能生效，架构状态改变 \\ \hline
\end{osgridlongtable}

\begin{workedexample}
执行 \code{ADD R0, R1}，初始 \(R0=5\)、\(R1=3\)。在 \(T_2\) 前，R0
仍是 5；\(T_3\) 形成内部候选 8，但软件可见 R0 仍未改变；只有
\(T_4\) 写回后 R0 才为 8。

若 \(T_1\) 取指返回错误，处理器应报告异常而不是继续使用旧 IR；若执行阶段
发现权限错误，精确异常要求不提交 R0。提交点把“内部计算过”与“架构上已经
发生”分开。
\end{workedexample}

\topic{硬连线控制与微程序控制是两种实现策略}
\term{硬连线控制}{hardwired control}直接由组合逻辑和状态机产生控制信号，
延迟可低但复杂指令修改困难；\term{微程序控制}{microprogrammed control}
把复杂指令拆成内部微操作序列，由控制存储器选择下一微指令。

x86 实现常把复杂指令译为内部微操作，再由乱序核心调度。微码不是操作系统
固件，也不是主板 SPI Flash 中的普通启动代码；它属于处理器实现层，更新和
装载方式由平台与处理器定义。

\topic{流水线让多条指令重叠，不让依赖自动消失}
若取指、译码、执行、访存、写回各需一拍，非流水实现一条指令需 5 拍；流水
稳定后理想情况下可每拍完成一条，但单条延迟仍约 5 拍。

\begin{center}
\begin{tikzpicture}[font=\scriptsize\sffamily,x=1.25cm,y=0.65cm]
  \foreach \x in {1,...,9} {
    \node at (\x,4.8) {拍 \x};
    \draw[BookRule] (\x-0.45,0) -- (\x-0.45,4.5);
  }
  \foreach \row/\name in {4/I1,3/I2,2/I3,1/I4} {
    \node[left] at (0.4,\row) {\name};
  }
  \foreach \x/\label/\colorname in {
    1/F/BookBlue,2/D/BookTeal,3/E/BookAmber,4/M/BookMuted,5/W/BookRed}
    \node[draw=\colorname,fill=\colorname!12,minimum width=9mm,minimum height=5mm]
      at (\x,4) {\label};
  \foreach \x/\label/\colorname in {
    2/F/BookBlue,3/D/BookTeal,4/E/BookAmber,5/M/BookMuted,6/W/BookRed}
    \node[draw=\colorname,fill=\colorname!12,minimum width=9mm,minimum height=5mm]
      at (\x,3) {\label};
  \foreach \x/\label/\colorname in {
    3/F/BookBlue,4/D/BookTeal,5/E/BookAmber,6/M/BookMuted,7/W/BookRed}
    \node[draw=\colorname,fill=\colorname!12,minimum width=9mm,minimum height=5mm]
      at (\x,2) {\label};
  \foreach \x/\label/\colorname in {
    4/F/BookBlue,5/D/BookTeal,6/E/BookAmber,7/M/BookMuted,8/W/BookRed}
    \node[draw=\colorname,fill=\colorname!12,minimum width=9mm,minimum height=5mm]
      at (\x,1) {\label};
  \node[align=left] at (10.2,2.5)
    {F 取指\\D 译码\\E 执行\\M 访存\\W 写回};
\end{tikzpicture}
\end{center}

流水线要处理：

\begin{itemize}[leftmargin=2.2em]
  \item 数据冒险：后一指令需要前一指令尚未写回的结果；
  \item 控制冒险：分支结果未定时不知道下一取指地址；
  \item 结构冒险：两阶段同时争用一个端口；
  \item 异常：内部已有多条在途指令，却要提供精确架构状态。
\end{itemize}

转发、停顿、预测、重命名和重排序缓冲分别缓解这些问题。它们属于微架构，
不能改变 ISA 对最终寄存器、内存顺序和异常的承诺。

\begin{failurebox}
框图中的一根“数据箭头”可能代表 8、64、256 bit 或更宽内部总线；一个方框
可能含数百万晶体管。框图用来证明职责和流向，不证明门级实现、物理位置或
时钟周期数。读图必须先确认抽象层级。
\end{failurebox}
